\chapter{Scalar response and stability}\label{sec:sc}\label{sec:scalar-stability}

In this section, we study the scalar Parisi equation of a single spin
under an arbitrary profile and the scalar response operator
$J^{\scal}_{a,d}$ built from it. The cavity interpolation of
\cref{sec:c7} decouples one spin from the others. At the independent
endpoint of that interpolation, the decoupled spin follows the scalar
law of \cref{sec:sc-law} with coefficient $d=\beta^2(1-t/N)$, and in
\cref{sec:c7-endpoint}, we express the derivative of the cavity covariance
at this endpoint through $J^{\scal}_{a,d}$, applied to the
covariance of the rest self-overlap. The contraction argument of
\cref{sec:c7-closure} needs the order bound
$J^{\scal}_{a,d}\preceq(1+\varepsilon_N)I$ with
$\varepsilon_N\le\rho/4$ at a minimizing profile. The goal of this
section is to prove this bound with
$\varepsilon_N\le C_{\rm st}[\rho^{-3/4}\eta^{1/4}+\lambda^2+N^{-1}]$,
where $\eta$ is a trial budget (\cref{prop:lz-stability}).

The argument has four steps.
\begin{enumerate}
\item We show in \cref{lem:sc-response} that, for every profile,
 $J_{a,d}$ is the derivative of the mean square magnetization
 $\Gamma_{a,d}$ when the covariance levels are transported by
 $q\mapsto q+\theta\varphi(q)$.
\item The scalar Parisi functional is convex. We compare its second
 derivative along a transport of levels with the lower bound given by
 convexity and the first variation, and we find that $J^{\scal}_{a,d}$
 is self-adjoint and that it is dominated by multiplication by the
 multiplier $k_{a,d}(q)=d\,\E^0h_q^2$, for every profile
 (\cref{lem:Jdom}).
\item At the scalar Parisi measure, we have $k=1$ on
 $\supp\mu_\beta=[0,q_\beta]$ and $k(q)\le e^{-\beta^2(q-q_\beta)/2}$
 above it (\cref{lem:reference}). On the support, this follows from an
 identity of Talagrand. Above the support, we use the log-concavity of
 a transformed density of the scalar diffusion and its heat evolution,
 both from \cite{Lopatto}.
\item The multiplier depends Lipschitz continuously on the profile in
 $L^1$ and on $d$ (\cref{lem:scalarcontinuity}). By
 \cref{prop:loc-isolation}, a finite-dimensional minimizer is close in
 $L^1$ to the scalar trial profile, which is close to
 $\beta\alpha_\beta$. Combining these facts with the first three steps,
 we obtain \cref{prop:lz-stability}.
\end{enumerate}
The bound is marginal: at the Parisi measure, $k=1$ on the support,
where the multiplier gives no contraction. The contraction in
\cref{sec:c7-closure} comes instead from the factor $t=1-\rho$, and it
needs only $\varepsilon_N\le\rho/4$. In the physics literature, De
Dominicis and Kondor~\cite[pp.~606--608]{DeDominicisKondor1983} found
that Parisi's solution with replica symmetry breaking is marginally stable
in a truncated replica model near the critical temperature. The
scalar identity on the support that
expresses this marginality appears in
\cite[Eqs.~(22)--(23)]{CrisantiRizzo2002}.

Sommers and Dupont gave an earlier variational formulation of the
coupled equations for the local magnetization and the distribution
of frozen fields~\cite[Section~2, Eqs.~(1)--(7), and
Eq.~(33)]{SommersDupont1984}; see also
\cite[Section~III, (14)--(18), (22)--(23)]{CrisantiRizzo2002}.
Those equations use the physical field and the order-parameter function
$q(x)$. Here we use the covariance level $q$ and the integrated
formulation in terms of probability measures described below.

We use the following results from earlier sections: the generalized
law and \cref{lem:fnd-posterior} of \cref{sec:gauss} for the basic
properties of the scalar law, \cref{lem:fnd-path-heat} for the
differentiation in \cref{lem:sc-response}, and, from \cref{sec:lz}, the
trial budget \eqref{eq:loc-trialbudget}, the root growth
\eqref{eq:lz-gsroot}, the description \eqref{eq:loc-scalar-measure} of
$\mu_\beta$, and \cref{prop:loc-isolation}. The results of this section
are used later as follows. The scalar law, described by \eqref{eq:sc-pde} and
\eqref{eq:sc-joint-law}, and the kernel bounds of \cref{lem:cl-scalar}
enter the cavity expansion in \cref{sec:c7,sec:c7-endpoint}. The
operator \eqref{eq:scalarJexplicit} and its kernel form
\eqref{eq:sc-kernel-form} are used in \cref{lem:linear}, \cref{lem:Jdom}
in \cref{prop:c7-closure}, and \cref{prop:lz-stability} in the proof of
\cref{thm:main} in \cref{sec:b7}.

\paragraph{Standing assumptions and notation.}
Throughout this section, $\beta>1$ is fixed and
$d\in[\beta^2/2,\beta^2]$. A \emph{profile} is the distribution function
$a(q)=\mu([0,q])$ of a Borel probability measure $\mu$ on $[0,1]$; it is
nondecreasing and right-continuous, and $a(1)=1$. We write
$\mu=s\delta_0+\nu$ with $s\ge0$, where the measure $\nu\ge0$ may charge
$0$. The scalar law depends only on $a$. The decomposition enters only
through $\nu$, which appears in the response operator. A \emph{step
profile} is a profile with finitely many values. Integrals over
$[0,q]$ include an atom at $q$, and integrals over $(q,1]$ exclude it.
Unless stated otherwise, $C_\beta$ denotes a constant depending only on
$\beta$, whose value may change from line to line. It does not depend
on $d\in[\beta^2/2,\beta^2]$, on $s$, or on the profile, including the
positions and masses of its atoms. No profile is assumed to minimize a
functional until \cref{sec:sc-minimizer}.

\section{The scalar law}\label{sec:sc-law}

The scalar Parisi equation and its diffusion are
\begin{equation}\label{eq:sc-pde}
 \begin{aligned}
 \partial_q u&=-\frac d2\bigl(\partial_{xx}u+a(q)(\partial_xu)^2\bigr),
 &u(1,x)&=\log(2\cosh x),\\
 \dd Y_q&=d\,a(q)\partial_xu(q,Y_q)\dd q+\sqrt d\dd B_q,
 &Y_0&=0,
 \end{aligned}
\end{equation}
where $B$ is a standard Brownian motion. We write $u=u_{a,d}$ when the
parameters need to be shown. Subscripts $x$ denote spatial derivatives,
and $u_j=\partial_x^ju$. When $a=\alpha_\beta$ and
$d=\beta^2$, the value $u(0,0)$ is the first term of the scalar Parisi
functional \eqref{eq:intro-parisi} at its minimizer.

\paragraph{Solutions and regularity.}
For a step profile, \eqref{eq:sc-pde} has a classical solution on each
interval where $a$ is constant. For a general profile, $u$ is the
uniform limit of the solutions for step profiles whose measures
converge weakly to $\mu$ \cite[Section~2.1]{Lopatto}. By
\cite[Proposition~2.1]{Lopatto}, applied with $\beta^2$ replaced by $d$,
the function $u(q,\cdot)$ is even, its spatial derivatives of every
positive order are bounded and jointly continuous on $[0,1]\times\R$,
and the spatial derivatives of the approximating solutions converge
uniformly \cite[(2.2)]{Lopatto}. The terminal value in \cite{Lopatto}
is $\log\cosh x$; it differs from ours by the constant $\log2$, which
shifts $u$ by $\log2$ and does not change its spatial derivatives. The
distribution functions of the approximating measures converge to $a$
at every continuity point of $a$, that is, at all but countably many
points. Passing to the limit in the equations integrated in the level
variable, we find that
\begin{equation}\label{eq:sc-pde-integrated}
 u(q,x)-u(r,x)=\frac d2\int_q^r\bigl(u_{xx}+a\,u_x^2\bigr)(v,x)\dd v,
 \qquad 0\le q\le r\le1,\ x\in\R .
\end{equation}
In particular, $u$ is Lipschitz in $q$, since the integrand is bounded.
For the weak space--time interpretation of the Parisi equation for
general measures, see \cite[Theorem~4]{JagannathTobasco2016}.
At an atom of $\mu$, the solution need not be differentiable in $q$. We
read \eqref{eq:sc-pde} and every other parabolic equation in this
section in this integrated form, including the linearized equations
derived below, and we never differentiate $a$.

\paragraph{Linear backward equations.}
Several proofs below use linear equations of the form
\begin{equation}\label{eq:sc-linear}
 \partial_qw+\frac{c(q)}2\,w_{xx}+b\,w_x+V\,w=-F,
\end{equation}
where $c$ is measurable with values in $[\beta^2/4,2\beta^2]$, and
$b$, $V$, and $F$ are bounded functions of $(q,x)$ that are measurable
in $q$ and have bounded continuous spatial derivatives. We use
\eqref{eq:sc-linear} only for functions $w$ that are Lipschitz in $q$
and have bounded continuous spatial derivatives of every order. We
read it in integrated form, as in \eqref{eq:sc-pde-integrated}. Fix
$(q,x)$, and let $Z$ solve $\dd Z_r=b(r,Z_r)\dd r+\sqrt{c(r)}\dd B_r$
for $r\ge q$, with $Z_q=x$. Itô's formula applies to
$r\mapsto e^{\int_q^rV(v,Z_v)\dd v}w(r,Z_r)$, with the derivative of
$w$ in $q$ taken almost everywhere; to see this, one may mollify $w$ in
the level variable and pass to the limit, using that the derivatives
are bounded. Taking expectations, we obtain the Feynman--Kac
representation
\begin{equation}\label{eq:sc-backward}
 w(q,x)=\E\Bigl[e^{\int_q^1V(v,Z_v)\dd v}\,w(1,Z_1)
   +\int_q^1e^{\int_q^rV(v,Z_v)\dd v}F(r,Z_r)\dd r\Bigr].
\end{equation}
In particular, if $w(1,\cdot)=0$, then
$\sup_x|w(q,x)|\le e^{\|V\|_\infty}\int_q^1\sup_x|F(r,x)|\dd r$.

\paragraph{Notation.}
We write $\Pp^0$ and $\E^0$ for probability and expectation under the
law of $Y$, together with a terminal spin $\epsilon\in\{-1,1\}$, which
is sampled given $Y_1$ with
\[
 \Pp^0(\epsilon=\pm1\mid Y_1)=\frac{e^{\pm Y_1}}{2\cosh Y_1}
 =e^{\pm Y_1-u(1,Y_1)}.
\]
Let $\calG^0_q$ be the filtration of $Y$; the terminal spin is not
$\calG^0_q$-measurable. We define
\[
 m_q=u_x(q,Y_q),\qquad h_q=u_{xx}(q,Y_q),\qquad
 \Gamma_{a,d}(q)=\E^0m_q^2,\qquad k_{a,d}(q)=d\,\E^0h_q^2 .
\]
We call $k_{a,d}$ the \emph{scalar multiplier}. For $0\le r\le\ell\le1$,
let $P^0_{r,\ell}\psi(y)=\E^0[\psi(Y_\ell)\mid Y_r=y]$ be the transition
operator of the Markov process $Y$, and set
\[
 f_p=m_p^2,\qquad \bar f_p=f_p-\E^0f_p,\qquad
 g_r^\ell(y)=P^0_{r,\ell}\bigl(u_x(\ell,\cdot)^2\bigr)(y),
\]
so that $g_r^\ell(Y_r)=\E^0[m_\ell^2\mid\calG_r^0]$. When we need the
diffusion started from $y$ at level $r$, we write $Y^{r,y}$ and
$\E_{r,y}$.

\begin{lemma}[Basic properties of the scalar law]\label{lem:sc-basic}
Let $a$ be a profile and $d\in[\beta^2/2,\beta^2]$.
\begin{enumerate}
\item[(a)] On $[0,1]\times\R$, we have
\begin{equation}\label{eq:sc-basic-bounds}
 |\partial_xu|\le1,\qquad
 0\le\partial_{xx}u\le1-(\partial_xu)^2\le1 .
\end{equation}
\item[(b)] For every $j\ge1$, there exists a constant $C_j(\beta)$,
depending only on $\beta$ and $j$, such that $|u_j|\le C_j(\beta)$ on
$[0,1]\times\R$. Further,
$\sup_x|u_{xx}(q,x)-u_{xx}(r,x)|\le C_\beta|q-r|$ for all
$q,r\in[0,1]$.
\item[(c)] We have
\begin{equation}\label{eq:sc-martingale}
 m_q=\E^0[\epsilon\mid\calG_q^0],\qquad m_0=0,\qquad
 \dd m_q=\sqrt d\,h_q\dd B_q,\qquad \Gamma_{a,d}'(q)=k_{a,d}(q).
\end{equation}
The function $k_{a,d}$ is continuous on $[0,1]$, and $\Gamma_{a,d}$ is
continuously differentiable on $[0,1]$, with one-sided derivatives at
the endpoints.
\end{enumerate}
\end{lemma}

\begin{proof}
We first identify the scalar law as a generalized law. We consider the
generalized law of \cref{def:gauss-law} with $N=1$, $\Lambda=d$, all
coupling coefficients equal to zero, root mass $s=0$, and the
probability measure $\mu$ in place of $\nu$; this choice is allowed for
every probability measure. The terminal value in
\eqref{eq:fnd-general-pde} is then
$\log\sum_{\sigma=\pm1}e^{x\sigma}=\log(2\cosh x)$, the Gaussian
disorder does not enter, and the outer density is one. Hence the
generalized law is the law of $(Y,\epsilon)$; it depends only on $a$.
For one spin, the conditional covariance is $C_q=1-m_q^2$, and
$H_q=h_q$.
\Cref{lem:fnd-posterior} implies that $m_q=\E^0[\epsilon\mid\calG^0_q]$,
that $m_0=0$, and that, for every $q\in[0,1]$, almost surely,
\begin{equation}\label{eq:sc-posterior-as}
 a(q)\bigl(1-m_q^2\bigr)\le h_q\le1-m_q^2 .
\end{equation}
Differentiating \eqref{eq:sc-pde-integrated} in $x$, we see that $u_x$
satisfies the integrated form of
$\partial_qu_x=-(d/2)(u_{xxx}+2a\,u_xu_{xx})$. Using Itô's formula
along $Y$, we find that the drift terms cancel and that
$\dd m_q=\sqrt d\,h_q\dd B_q$. This is the case $N=1$ of
\eqref{eq:gauss-martingale}.

\emph{Proof of (a).} Fix $q\in(0,1]$. Since the drift of $Y$ is bounded
by $d$, Girsanov's theorem implies that the law of $Y_q$ is equivalent
to the centered Gaussian law with variance $dq$. In particular, it
charges every nonempty open interval. The functions
\[
 u_{xx}(q,\cdot)-a(q)\bigl(1-u_x(q,\cdot)^2\bigr),\qquad
 1-u_x(q,\cdot)^2-u_{xx}(q,\cdot),\qquad
 1-u_x(q,\cdot)^2
\]
are continuous, and by \eqref{eq:sc-posterior-as} and $|m_q|\le1$, each
of them is nonnegative at $Y_q$ almost surely. A continuous function
with this property is nonnegative at every $x$. This proves \eqref{eq:sc-basic-bounds} for
$q>0$, since $a(q)(1-u_x^2)\ge0$. Letting $q\downarrow0$ and using the
joint continuity of $u_x$ and $u_{xx}$, we see that
\eqref{eq:sc-basic-bounds} also holds at $q=0$. (With the right continuity of
$a$, the lower bound $u_{xx}\ge a(q)(1-u_x^2)$ also holds everywhere.)

\emph{Proof of (b).} We differentiate the equation and use the
Feynman--Kac representation, following the method of
\cite[Lemma~16]{JagannathTobasco2016}. By \eqref{eq:sc-basic-bounds}, we have
$|u_1|\le1$ and $|u_2|\le1$. For $j\ge3$, the Leibniz rule gives
\[
 \partial_x^j(u_x^2)=2u_1u_{j+1}+2j\,u_2u_j
 +\sum_{i=2}^{j-2}\binom ji u_{i+1}u_{j-i+1},
\]
with an empty sum for $j=3$. Differentiating
\eqref{eq:sc-pde-integrated} $j$ times in $x$, we see that $u_j$ solves
\eqref{eq:sc-linear} with
\[
 c=d,\qquad b=d\,a\,u_1,\qquad V=j\,d\,a\,u_2,\qquad
 F=\frac{d\,a}2\sum_{i=2}^{j-2}\binom ji u_{i+1}u_{j-i+1},
\]
and terminal value $\partial_x^j\log(2\cosh x)$. Recall that the
derivatives of $u$ are bounded by \cite[Proposition~2.1]{Lopatto}; in
particular, the hypotheses of \eqref{eq:sc-backward} hold, with $Z=Y$.
By \eqref{eq:sc-basic-bounds}, $V$ takes values in $[0,j\beta^2]$; note
that the exponent in \eqref{eq:sc-backward} is nonnegative. Then
\eqref{eq:sc-backward} implies that
\[
 \|u_j\|_\infty\le e^{j\beta^2}\Bigl(
 \|\partial_x^j\log(2\cosh x)\|_\infty
 +\frac{\beta^2}2\sum_{i=2}^{j-2}\binom ji
          \|u_{i+1}\|_\infty\|u_{j-i+1}\|_\infty\Bigr).
\]
Every derivative on the right has order between $3$ and $j-1$. By
induction on $j$, we conclude that $\|u_j\|_\infty\le C_j(\beta)$ for
all $j\ge1$, uniformly over profiles and $d\in[\beta^2/2,\beta^2]$.
Differentiating \eqref{eq:sc-pde-integrated} twice in $x$, we find
\[
 u_{xx}(q,x)-u_{xx}(r,x)=\frac d2\int_q^r
 \Bigl[u_4+2a\bigl(u_2^2+u_1u_3\bigr)\Bigr](v,x)\dd v,
\]
and the integrand is bounded by a constant depending only on $\beta$.
This proves the Lipschitz bound in (b).

\emph{Proof of (c).} The first three identities in
\eqref{eq:sc-martingale} were shown above. By Itô's isometry and
$m_0=0$, we have
\[
 \Gamma_{a,d}(q)=\E^0m_q^2=d\int_0^q\E^0h_r^2\dd r
 =\int_0^qk_{a,d}(r)\dd r .
\]
By (b), we have $|h_q-h_r|\le C_\beta|q-r|+C_3(\beta)|Y_q-Y_r|$. Since
$Y$ has continuous paths and $|h|\le1$, dominated convergence implies
that $k_{a,d}$ is continuous. It follows that $\Gamma_{a,d}$ is
continuously differentiable on $[0,1]$, with $\Gamma_{a,d}'=k_{a,d}$.
\end{proof}

We also use the exact path likelihood of the scalar law, for every
profile, including profiles with atoms. It is the change of measure of
\cite[Lemma~8.5 and Corollary~8.7]{JagannathTobasco2017PhaseDiagram};
we include the derivation in order to record the conventions at atoms.
It extends the likelihood used for step profiles in the proof of
\cref{lem:fnd-posterior} to every profile, without approximation.

\begin{lemma}[Exact path likelihood]\label{lem:sc-likelihood}
Let $a$ be a profile and $d\in[\beta^2/2,\beta^2]$. Under a reference
law, let $X_q=\sqrt d\,B_q$, where $B$ is a standard Brownian motion.
For every $q\in[0,1]$, the law of $(Y_r)_{0\le r\le q}$ has density
\begin{equation}\label{eq:sc-prefix-law}
 \exp\Bigl\{a(q)u(q,X_q)-\int_{[0,q]}u(r,X_r)\,\mu(\dd r)\Bigr\}
\end{equation}
with respect to the law of $(X_r)_{0\le r\le q}$. The joint law of
$(Y,\epsilon)$ has density
\begin{equation}\label{eq:sc-joint-law}
 \exp\Bigl\{\epsilon X_1-su(0,0)-\int_{[0,1]}u(q,X_q)\,\nu(\dd q)\Bigr\}
\end{equation}
with respect to the product of the law of $X$ and counting measure on
$\{-1,1\}$. The same holds if the coefficient $d$ in \eqref{eq:sc-pde}
is replaced by a bounded measurable function $c(q)>0$ almost
everywhere, and $X_q=\int_0^q\sqrt{c(r)}\dd B_r$.
\end{lemma}

\begin{proof}
We write $\Upsilon_r=u(r,X_r)$. The drift $d\,a\,u_x$ of $Y$ is
bounded, and by \eqref{eq:sc-basic-bounds}, we have
\[
 \exp\Bigl\{\frac d2\int_0^qa(r)^2u_x(r,X_r)^2\dd r\Bigr\}\le e^{\beta^2/2},
\]
which is Novikov's condition. By Girsanov's theorem, the density of the law of $(Y_r)_{r\le q}$ with
respect to the law of $(X_r)_{r\le q}$ is
\[
 \exp\Bigl\{\int_0^qa(r)u_x(r,X_r)\dd X_r
 -\frac d2\int_0^qa(r)^2u_x(r,X_r)^2\dd r\Bigr\}.
\]
Under the reference law, using Itô's formula and
\eqref{eq:sc-pde-integrated}, we have
\[
 \dd\Upsilon_r=u_x(r,X_r)\dd X_r-\frac d2a(r)u_x(r,X_r)^2\dd r .
\]
The exponent therefore equals $\int_0^qa(r)\dd\Upsilon_r$. The function $a$ is
right-continuous and of bounded variation, with $\dd a=\mu$ on
$(0,1]$, and $\Upsilon$ is a continuous semimartingale. Integrating by
parts, we have
\[
 a(q)\Upsilon_q-a(0)\Upsilon_0
 =\int_{(0,q]}a(r-)\dd\Upsilon_r+\int_{(0,q]}\Upsilon_r\,\mu(\dd r).
\]
The functions $a(r-)$ and $a(r)$ differ only at the countably many
atoms of $\mu$, which affects neither the stochastic integral nor the
Lebesgue integral in $\dd\Upsilon_r$. Since $X_0=0$ and
$a(0)=\mu(\{0\})$, we have
$a(0)\Upsilon_0=\int_{\{0\}}u(r,X_r)\,\mu(\dd r)$. Inserting these facts
into the previous display, we obtain
\[
 \int_0^qa(r)\dd\Upsilon_r=a(q)u(q,X_q)-\int_{[0,q]}u(r,X_r)\,\mu(\dd r).
\]
This proves \eqref{eq:sc-prefix-law}. At an atom of $\mu$ at $q$, the
jump of the first term is canceled by the same atom in the integral,
and at $q=0$ the density equals one.

Given the path, the terminal spin has probability
$e^{\epsilon Y_1-u(1,Y_1)}$ relative to counting measure. Multiplying
\eqref{eq:sc-prefix-law} at $q=1$ by this factor, and using $a(1)=1$
and $\mu=s\delta_0+\nu$, we obtain \eqref{eq:sc-joint-law}. No further
normalization is needed, since both factors are probability densities.

For a coefficient $c(q)$ that is positive almost everywhere, we first
choose the everywhere positive representative
$c^+(q)=c(q)$ when $c(q)>0$ and $c^+(q)=1$ otherwise. It is measurable
and bounded and equals $c$ almost everywhere. This change preserves
every prefix clock, the integrated equation and drift, and the
stochastic integral defining $X$. Thus the law and the likelihood
are independent of this choice of representative.
Let $C(q)=\int_0^qc(r)\dd r$. The substitution
$\tau=C(q)/C(1)$ turns the equation with coefficient $c(q)$ into
\eqref{eq:sc-pde} with coefficient $C(1)$ and the profile of the image
of $\mu$ under $q\mapsto C(q)/C(1)$. This substitution is made in the
integrated equation, using $\dd\tau=c(q)\dd q/C(1)$. Since $c>0$ almost
everywhere, $q\mapsto C(q)/C(1)$ is an increasing homeomorphism of
$[0,1]$; no positive lower bound on $c$ is required. By
\cite[Proposition~2.1]{Lopatto}, applied with coefficient $C(1)$, the
solution satisfies $|u_x|\le1$ and has bounded spatial
derivatives. The argument above then applies with $d\dd r$ replaced by
$c(r)\dd r$; in Novikov's condition, $\beta^2$ is replaced by an
upper bound for $c$.
\end{proof}

\section{Continuity and kernel bounds}
\label{sec:sc-kernels}

The proof of the next lemma follows the comparison argument of
\cite[Lemma~14]{JagannathTobasco2016}, which bounds the differences of
the solutions and of their first spatial derivatives by the $L^1$
distance of the profiles. Here we also allow $d$ to vary and compare
the second derivatives along the diffusion.

\begin{lemma}[Continuity of the scalar multiplier]\label{lem:scalarcontinuity}
There exists $C_\beta<\infty$ such that, for any two profiles
$a,\widetilde a$ and any $d,\widetilde d\in[\beta^2/2,\beta^2]$,
\begin{equation}\label{eq:scalarcontinuity}
 \sup_{0\le q\le1}|k_{a,d}(q)-k_{\widetilde a,\widetilde d}(q)|
 \le C_\beta\bigl(\|a-\widetilde a\|_{L^1(0,1)}+|d-\widetilde d|\bigr).
\end{equation}
\end{lemma}

\begin{proof}
Let $u$ and $\widetilde u$ solve \eqref{eq:sc-pde} for the two
parameter pairs, and let $Y$ and $\widetilde Y$ be the corresponding
diffusions. We set $\hat u=u-\widetilde u$ and
$\Delta=\|a-\widetilde a\|_{L^1(0,1)}+|d-\widetilde d|$. The proof has
three steps. We first bound $\hat u_x$ and $\hat u_{xx}$ by
$C_\beta\Delta$, then couple the two diffusions, and finally compare
the multipliers.

\paragraph{Step 1: the difference of the solutions.}
Subtracting the two equations and writing
$u_x^2-\widetilde u_x^2=(u_x+\widetilde u_x)\hat u_x$, we find that
$\hat u$ solves \eqref{eq:sc-linear} with
\[
 c=d,\qquad b=\omega,\qquad V=0,\qquad F=\Xi,\qquad
 \hat u(1,\cdot)=0,
\]
where
\[
 \omega=\frac{d\,a(q)}2\bigl(u_x+\widetilde u_x\bigr),\qquad
 \Xi=\frac{d-\widetilde d}2\,\widetilde u_{xx}
     +\delta(q)\,\widetilde u_x^2,\qquad
 \delta(q)=\frac{d\,a(q)-\widetilde d\,\widetilde a(q)}2 .
\]
Using \cref{lem:sc-basic}(a) and~(b), we see that $0\le\omega_x\le\beta^2$ and
$|\omega_{xx}|\le C_\beta$. Further,
\[
 \Xi_x=\frac{d-\widetilde d}2\,\widetilde u_3
       +2\delta\,\widetilde u_1\widetilde u_2,\qquad
 \Xi_{xx}=\frac{d-\widetilde d}2\,\widetilde u_4
       +2\delta\bigl(\widetilde u_2^2+\widetilde u_1\widetilde u_3\bigr),
\]
and $2|\delta(q)|\le|d-\widetilde d|+\beta^2|a(q)-\widetilde a(q)|$.
Using \cref{lem:sc-basic}(b) again, we find
\[
 \int_0^1\sup_x\bigl(|\Xi_x|+|\Xi_{xx}|\bigr)(q,x)\dd q
 \le C_\beta\int_0^1\bigl(|d-\widetilde d|+|a(q)-\widetilde a(q)|\bigr)\dd q
 \le C_\beta\Delta .
\]
Differentiating the equation for $\hat u$ once and twice in $x$, we see
that $w_1=\hat u_x$ and $w_2=\hat u_{xx}$ solve \eqref{eq:sc-linear}
with $c=d$, $b=\omega$, zero terminal values, and
\[
 V=\omega_x,\quad F=\Xi_x\qquad\text{and}\qquad
 V=2\omega_x,\quad F=\Xi_{xx}+\omega_{xx}w_1,
\]
respectively. Then \eqref{eq:sc-backward} implies that
\[
 \begin{aligned}
 \|w_1\|_\infty&\le e^{\beta^2}\int_0^1\|\Xi_x(q,\cdot)\|_\infty\dd q
 \le C_\beta\Delta,\\
 \|w_2\|_\infty&\le e^{2\beta^2}\int_0^1\bigl(\|\Xi_{xx}(q,\cdot)\|_\infty
 +\|\omega_{xx}\|_\infty\|w_1\|_\infty\bigr)\dd q\le C_\beta\Delta .
 \end{aligned}
\]
This shows that
\begin{equation}\label{eq:sc-pde-distance}
 \sup_{q,x}\bigl(|u_x-\widetilde u_x|+|u_{xx}-\widetilde u_{xx}|\bigr)
 \le C_\beta\Delta .
\end{equation}

\paragraph{Step 2: coupling the diffusions.}
We drive $Y$ and $\widetilde Y$ by the same Brownian motion and set
$D(q)=\sup_{r\le q}|Y_r-\widetilde Y_r|$. By
\eqref{eq:sc-basic-bounds}, the drift $d\,a(r)u_x(r,\cdot)$ is
$\beta^2$-Lipschitz in space, and at a common point it differs from
$\widetilde d\,\widetilde a(r)\widetilde u_x(r,\cdot)$ by at most
$|d-\widetilde d|+\beta^2|a(r)-\widetilde a(r)|+\beta^2C_\beta\Delta$,
where we used \eqref{eq:sc-pde-distance}. Subtracting the integral
equations for $Y$ and $\widetilde Y$, we obtain
\[
 D(q)\le\beta^2\int_0^qD(r)\dd r+C_\beta\Delta
 +|\sqrt d-\sqrt{\widetilde d}\,|\sup_{r\le1}|B_r| .
\]
Since $d,\widetilde d\ge\beta^2/2$, we have
$|\sqrt d-\sqrt{\widetilde d}\,|\le|d-\widetilde d|/(\sqrt2\,\beta)$.
By Gronwall's inequality and $\E\sup_{r\le1}|B_r|\le2$, we conclude
that
\begin{equation}\label{eq:sc-coupling}
 \E\sup_{q\le1}|Y_q-\widetilde Y_q|
 \le C_\beta\bigl(\Delta+|\sqrt d-\sqrt{\widetilde d}\,|\,
                     \E\sup_{q\le1}|B_q|\bigr)\le C_\beta\Delta .
\end{equation}
The same argument, applied to the two diffusions started from a common
point $y$ at a common level $r$, yields the same bound, uniformly in
$(r,y)$.

\paragraph{Step 3: the multipliers.}
We write $h=u_{xx}(q,Y_q)$ and $\widetilde h=\widetilde u_{xx}(q,\widetilde Y_q)$.
Note that $|h^2-\widetilde h^2|\le2|h-\widetilde h|$, because
$0\le h,\widetilde h\le1$. Then
\[
 \begin{aligned}
 |k_{a,d}(q)-k_{\widetilde a,\widetilde d}(q)|
 &\le|d-\widetilde d|+2\beta^2\,\E|h-\widetilde h|\\
 &\le|d-\widetilde d|+2\beta^2\Bigl(\|u_{xx}-\widetilde u_{xx}\|_\infty
   +\|u_3\|_\infty\,\E|Y_q-\widetilde Y_q|\Bigr).
 \end{aligned}
\]
Inserting \eqref{eq:sc-pde-distance}, \eqref{eq:sc-coupling}, and the
bound $\|u_3\|_\infty\le C_3(\beta)$ into the last line, we obtain
\eqref{eq:scalarcontinuity}.
\end{proof}

The next lemma collects the bounds on the kernels of the response
operator that the cavity estimates of \cref{sec:c7-endpoint} use. Each
bound retains a factor $\Gamma$ at the relevant levels.

\begin{lemma}[Scalar kernels and their level weight]\label{lem:cl-scalar}
For every profile $a$ and $d\in[\beta^2/2,\beta^2]$, write
$\Gamma=\Gamma_{a,d}$. Then $\Gamma$ is nondecreasing and
\begin{equation}\label{eq:cl-Gamma-bound}
 0\le\Gamma(p)\le\min\{1,dp\}\le\min\{1,\beta^2p\}.
\end{equation}
For all $p,r,\ell\in[0,1]$,
\begin{align}
 |\E^0[m_ph_pm_r]|&\le\sqrt{\Gamma(p)\Gamma(r)},&
 |\E^0[\bar f_p\epsilon m_r]|&\le\sqrt{\Gamma(p)\Gamma(r)},
                                                        \label{eq:cl-k1}\\
 |\E^0[\bar f_pm_rm_\ell]|&\le\sqrt{\Gamma(p)\Gamma(r\wedge\ell)},&
 |\E^0[\bar f_pm_pm_r]|&\le\sqrt{\Gamma(p)\Gamma(p\wedge r)},
                                                        \label{eq:cl-k2}\\
 0\le\Var^0(f_p)&\le\Gamma(p),&
 0\le\E^0[m_p^2h_p]&\le\Gamma(p).                       \label{eq:cl-k3}
\end{align}
For $\ell\ge p$,
\begin{equation}\label{eq:cl-k4}
 |\Cov^0(f_p,g_p^\ell(Y_p))|\le\sqrt{\Gamma(p)\Gamma(\ell)}.
\end{equation}
\end{lemma}

\begin{proof}
By \cref{lem:sc-basic}(c) and \eqref{eq:sc-basic-bounds}, we have
$\Gamma(p)=d\int_0^p\E^0h_r^2\dd r$, where $0\le h_r\le1$. This implies
that $\Gamma$ is nondecreasing and that $0\le\Gamma(p)\le dp$. We also
have $\Gamma\le1$, because $|m_p|\le1$. This proves \eqref{eq:cl-Gamma-bound}. Using $|m_p|\le1$
again, we have
\[
 \Var^0(f_p)=\E^0\bar f_p^2\le\E^0m_p^4\le\E^0m_p^2=\Gamma(p).
\]
The bounds below follow from the Cauchy--Schwarz inequality, together
with $|h_p|\le1$, $|\epsilon|=1$, and $|m|\le1$:
\[
 \begin{aligned}
 |\E^0[m_ph_pm_r]|^2
 &\le\E^0[m_p^2h_p^2]\,\E^0m_r^2\le\Gamma(p)\Gamma(r),\\
 |\E^0[\bar f_p\epsilon m_r]|^2
 &\le\E^0\bar f_p^2\,\E^0m_r^2\le\Gamma(p)\Gamma(r),\\
 |\E^0[\bar f_pm_rm_\ell]|^2
 &\le\E^0\bar f_p^2\,\E^0[m_r^2m_\ell^2]\le\Gamma(p)\Gamma(r\wedge\ell).
 \end{aligned}
\]
In the last line, we used $m_r^2m_\ell^2\le m_{r\wedge\ell}^2$, which
holds because the factor at the smaller level is one of the two factors
and the other is at most one. This proves \eqref{eq:cl-k1} and the first
bound in \eqref{eq:cl-k2}; the second bound in \eqref{eq:cl-k2} is the
case $\ell=p$. The bounds \eqref{eq:cl-k3} follow from the estimate on
$\Var^0(f_p)$ and from $0\le h_p\le1$, which gives
$0\le\E^0[m_p^2h_p]\le\E^0m_p^2$. Finally, let $\ell\ge p$. We have
$0\le g_p^\ell\le1$ and $\E^0g_p^\ell(Y_p)=\E^0m_\ell^2=\Gamma(\ell)$,
which imply that
\[
 \Var^0\bigl(g_p^\ell(Y_p)\bigr)\le\E^0\bigl(g_p^\ell(Y_p)\bigr)^2
 \le\E^0g_p^\ell(Y_p)=\Gamma(\ell),
\]
and \eqref{eq:cl-k4} follows from the Cauchy--Schwarz inequality and
the bound on $\Var^0(f_p)$.
\end{proof}

\section{The response operator}\label{sec:sc-response}

For every profile, the operator defined below is the response of
$\Gamma$ to a transport of covariance levels (\cref{lem:sc-response}).
In \cite[Lemma~5.2, (H2)]{OpenAILimitingLaw}, a scalar response is
written as a diagonal multiplier minus a positive semidefinite Hessian
kernel. Here we identify the second variation of the scalar Parisi
functional along transports of levels, and we compare it using the
convexity of this functional along segments of distribution functions
\cite{AuffingerChen2015Uniqueness}. In \cref{sec:sc-domination}, this
shows that the scalar response operator is dominated by multiplication by
$k_{a,d}$ at every profile. The transports fix the levels $0$ and $1$,
and the operator acts on the \emph{scalar transport space}
\[
 \mathcal H_\nu=\bigl\{v\in L^2(\nu):
                 v(0)=0\text{ if }\nu(\{0\})>0,
                 \ v(1)=0\text{ if }\nu(\{1\})>0\bigr\}.
\]
This is a closed subspace of $L^2(\nu)$, and the orthogonal projection
onto it is multiplication by $\ind_{(0,1)}$.

The operator involves the derivative of $g_r^\ell$ in its argument.
Since the noise of $Y$ is additive, the diffusion $Y^{r,y}$ started
from $y$ at level $r$ solves an ordinary differential equation driven
by the Brownian path, and it is differentiable in $y$. Its derivative
is $\xi_\ell=\exp\{d\int_r^\ell a(v)u_{xx}(v,Y^{r,y}_v)\dd v\}$, which
lies in $[1,e^{\beta^2}]$ by \eqref{eq:sc-basic-bounds}.
Differentiating $g_r^\ell(y)=\E[u_x(\ell,Y^{r,y}_\ell)^2]$ under the
expectation, we find
\begin{equation}\label{eq:sc-dg}
 \partial_yg_r^\ell(y)=\E\bigl[2u_x(\ell,Y_\ell^{r,y})
 u_{xx}(\ell,Y_\ell^{r,y})\,\xi_\ell\bigr],\qquad
 |\partial_yg_r^\ell|\le2e^{\beta^2}.
\end{equation}

\begin{definition}[Response operator]\label{def:sc-response}
Let $a$ be a profile with $\mu=s\delta_0+\nu$, and let
$d\in[\beta^2/2,\beta^2]$. For a bounded Borel function $\varphi$ on
$[0,1]$ and $p\in[0,1]$, we define
\begin{align}
 \frac1d(J_{a,d}\varphi)(p)={}&
 \E^0\bigl[h_p^2+2(1-a(p))m_p^2h_p\bigr]\varphi(p)
 -\int_{(p,1]}\E^0\bigl[m_p\partial_y g_p^\ell(Y_p)\bigr]
                                      \varphi(\ell)\,\nu(\dd\ell)\nonumber\\
 &-2\int\E^0[m_ph_pm_r]\varphi(p\wedge r)\,\nu(\dd r)
  -\int\E^0[\bar f_p\epsilon m_r]\varphi(r)\,\nu(\dd r)\nonumber\\
 &+\frac12\int a(r)\E^0[\bar f_pm_r^2]\varphi(r)\,\nu(\dd r)
  +\frac12\iint\E^0[\bar f_pm_rm_\ell]\varphi(r\wedge\ell)
                                      \,\nu(\dd r)\nu(\dd\ell)\nonumber\\
 &+\frac12\int\nu(\dd r)\int_{(r,1]}
           \E^0\bigl[\bar f_pg_r^\ell(Y_r)\bigr]\varphi(\ell)\,\nu(\dd\ell).
 \label{eq:scalarJexplicit}
\end{align}
Integrals without a domain are over $[0,1]$. The same formula defines
$J_{a,d}v$ for $v\in L^2(\nu)$. The \emph{scalar response operator}
$J^{\scal}_{a,d}$ is the compression of $J_{a,d}$ to $\mathcal H_\nu$,
that is, $J^{\scal}_{a,d}v=\ind_{(0,1)}J_{a,d}v$ for $v\in\mathcal H_\nu$.
\end{definition}

At an atom of $\nu$ at $p$, the right-continuous value $a(p)$ includes
the atom, and the strict future interval $(p,1]$ excludes it. For
$v\in L^2(\nu)$, the formula involves $v$ only through $v(p)$ and
through values at points that are integrated against $\nu$, because
$p\wedge r\in\{p,r\}$ and $r\wedge\ell\in\{r,\ell\}$. We show in the
next paragraph that it defines $J_{a,d}v$ as an element of $L^2(\nu)$.

\paragraph{Kernel form and boundedness.}
We now rewrite \eqref{eq:scalarJexplicit}. For $r>p$, conditioning on
$\calG^0_p$ and using \eqref{eq:sc-martingale}, we have
$\E^0[m_ph_pm_r]=\E^0[m_p^2h_p]$, and $\nu((p,1])=1-a(p)$. Splitting the
third term of \eqref{eq:scalarJexplicit} at $r=p$, we see that its part
over $(p,1]$ equals $-2(1-a(p))\E^0[m_p^2h_p]\varphi(p)$, which cancels
the second part of the first coefficient. In the double integral, we
split $\varphi(r\wedge\ell)$ into the regions $r<\ell$, $\ell<r$, and the
diagonal $r=\ell$, which carries mass only at atoms of $\nu$; the first
two regions give equal contributions by symmetry. In the last term, we
exchange the order of integration. All integrands are bounded, and
Fubini's theorem justifies these steps. The result is
\begin{equation}\label{eq:sc-kernel-form}
 (J_{a,d}v)(p)=k_{a,d}(p)v(p)+d\int\mathcal K_{a,d}(p,r)v(r)\,\nu(\dd r),
\end{equation}
for every bounded Borel $v$, where
\begin{equation}\label{eq:sc-kernel}
 \begin{aligned}
 \mathcal K_{a,d}(p,r)={}&
 -\ind_{\{r>p\}}\E^0\bigl[m_p\partial_yg_p^r(Y_p)\bigr]
 -2\,\ind_{\{r\le p\}}\E^0[m_ph_pm_r]
 -\E^0[\bar f_p\epsilon m_r]\\
 &+\frac12a(r)\E^0[\bar f_pm_r^2]
 +\int_{(r,1]}\E^0[\bar f_pm_rm_\ell]\,\nu(\dd\ell)
 +\frac12\nu(\{r\})\E^0[\bar f_pm_r^2]\\
 &+\frac12\int_{[0,r)}\E^0\bigl[\bar f_pg_\ell^r(Y_\ell)\bigr]\,\nu(\dd\ell).
 \end{aligned}
\end{equation}
The term with $\nu(\{r\})$ is the diagonal part of the double integral.
The kernel $\mathcal K_{a,d}$ is a bounded Borel function. Indeed, the
first two terms are never both nonzero and are bounded by
$2e^{\beta^2}$ by \eqref{eq:sc-dg}; since $|\bar f_p|\le1$, $|m|\le1$, and $0\le g\le1$, the
third and fourth terms are bounded by $1$ and $1/2$, and the last three
terms together by $\nu([0,1])\le1$. We conclude that
$|\mathcal K_{a,d}|\le2e^{\beta^2}+3$. By the Cauchy--Schwarz
inequality,
$\|\int\mathcal K_{a,d}(\cdot,r)v(r)\,\nu(\dd r)\|_{L^2(\nu)}
\le(2e^{\beta^2}+3)\nu([0,1])\|v\|_{L^2(\nu)}$. It follows that
\eqref{eq:sc-kernel-form} extends to $v\in L^2(\nu)$, where it agrees
with \eqref{eq:scalarJexplicit}. Using $0\le k_{a,d}\le d\le\beta^2$,
we find $\|J_{a,d}\|_{L^2(\nu)\to L^2(\nu)}\le\beta^2(4+2e^{\beta^2})$.
Compression does not increase the norm, which gives
\begin{equation}\label{eq:sc-J-bounded}
 \|J^{\scal}_{a,d}\|_{\mathcal H_\nu\to\mathcal H_\nu}\le C_\beta .
\end{equation}

\paragraph{Transports of levels.}
Let $\varphi$ be the restriction to $[0,1]$ of a smooth function, with
$\varphi(0)=\varphi(1)=0$, and define $T_\theta(q)=q+\theta\varphi(q)$.
We fix $\theta_0\in(0,1]$ such that $\theta_0\|\varphi'\|_\infty\le1/2$.
For $|\theta|\le\theta_0$, the map $T_\theta$ is an
increasing bijection of $[0,1]$ with $T_\theta'\in[1/2,3/2]$. We replace
the covariance $d\,(q\wedge r)$ of $X$ in \eqref{eq:sc-joint-law} by
$d\,T_\theta(q\wedge r)$ while keeping $\nu$. Equivalently, we replace
the coefficient $d$ in \eqref{eq:sc-pde} by
\[
 c_\theta(q)=d\,T_\theta'(q)=d\bigl(1+\theta\varphi'(q)\bigr)
 \in\Bigl[\frac d2,\frac{3d}2\Bigr]\subset\Bigl[\frac{\beta^2}4,2\beta^2\Bigr],
\]
and keep $a$. We write $u_\theta$ and $Y^\theta$ for the resulting
solution and diffusion, $u_{\theta,x}=\partial_xu_\theta$ and similarly
for higher derivatives, and
\[
 \widetilde\Gamma_\theta(p)=\E\bigl[\partial_xu_\theta(p,Y^\theta_p)^2\bigr]
\]
for the resulting mean square magnetization. Let
$a_\theta=a\circ T_\theta^{-1}$ be the distribution function of the
image measure $T_{\theta\#}\mu$. Under the substitution $x'=T_\theta(q)$
in the level variable, the equation for $u_\theta$ becomes
\eqref{eq:sc-pde} with coefficient $d$ and profile $a_\theta$, and the
diffusion $Y^\theta$ becomes a time change of the diffusion for
$a_\theta$. In formulas,
\begin{equation}\label{eq:sc-transported}
 u_\theta(q,x)=u_{a_\theta,d}(T_\theta(q),x),\qquad
 \widetilde\Gamma_\theta(p)=\Gamma_{a_\theta,d}(T_\theta(p)),
\end{equation}
that is, $(u_\theta,Y^\theta)$ is the scalar law of $T_{\theta\#}\mu$
with coefficient $d$, observed at the transported levels. In
particular, $\widetilde\Gamma_0=\Gamma_{a,d}$, and by
\cref{lem:sc-basic}(b), the spatial derivatives of $u_\theta$ of each
order are bounded uniformly in $|\theta|\le\theta_0$ and over profiles.

\begin{lemma}[Response identity]\label{lem:sc-response}
Let $a$ be a profile, let $d\in[\beta^2/2,\beta^2]$, and let $\varphi$
and $\widetilde\Gamma_\theta$ be as above. Then, for every $p\in[0,1]$,
the function $\theta\mapsto\widetilde\Gamma_\theta(p)$ is continuously
differentiable near $\theta=0$, and
\[
 \frac{\dd}{\dd\theta}\widetilde\Gamma_\theta(p)\Big|_{\theta=0}
 =(J_{a,d}\varphi)(p).
\]
\end{lemma}

\begin{proof}
The proof has four steps. In Step 1, we show that $u_\theta$ and its
spatial derivatives are continuously differentiable in $\theta$; the
argument works with the integrated equations and never differentiates
$a$. In Step 2, we compute the parameter derivative of $u_\theta$ and
prove \eqref{eq:scalarNodeVariation}. In Step 3, we write
$\widetilde\Gamma_\theta(p)$ as a normalized expectation on a fixed
Gaussian path space and differentiate it with
\cref{lem:fnd-path-heat}. In Step 4, we match the resulting terms with
those of \eqref{eq:scalarJexplicit}. Throughout the proof,
$|\theta|\le\theta_0$, and constants may depend on $\beta$ and
$\varphi$.

\paragraph{Step 1: differentiability in the parameter.}
Fix $\theta$, and for small $\delta\ne0$ let $v=u_{\theta+\delta}$,
$u=u_\theta$, and $D_\delta=(v-u)/\delta$. Since
$c_{\theta+\delta}-c_\theta=\delta d\varphi'$, subtracting the two
equations and dividing by $\delta$, we find
\begin{equation}\label{eq:sc-Dh}
 \partial_qD_\delta+\frac{c_\theta}2(D_\delta)_{xx}
       +\frac{c_\theta a}2(v_x+u_x)(D_\delta)_x
 =-\frac{d\varphi'}2\bigl(v_{xx}+a\,v_x^2\bigr),\qquad D_\delta(1,\cdot)=0 .
\end{equation}
This is \eqref{eq:sc-linear} with $c=c_\theta$, $V=0$, and a drift and
a source whose spatial derivatives of each order are bounded uniformly
in $\delta$ and $\theta$. Differentiating \eqref{eq:sc-Dh} $j$ times in
$x$, we see that $\partial_x^jD_\delta$ solves an equation of the same
form, with zeroth-order coefficient $j$ times the spatial derivative
of the drift, and with a source that involves only derivatives of
$D_\delta$ of order less than $j$ and spatial derivatives of the drift
and of the source of \eqref{eq:sc-Dh}. By \eqref{eq:sc-backward} and
induction on $j$, we have
$\sup_{q,x}|\partial_x^jD_\delta|\le C_j$ for every $j\ge0$, uniformly
in $\delta$ and $\theta$. In particular, $v-u=\delta D_\delta$ and all
its spatial derivatives are $O(|\delta|)$.

Next, we write the drift in \eqref{eq:sc-Dh} as
$c_\theta a\,u_x+(c_\theta a/2)(v_x-u_x)$ and move the second part
to the right side. For two increments $\delta$ and $\delta'$, the
difference $D_\delta-D_{\delta'}$ then solves \eqref{eq:sc-linear} with
drift $c_\theta a\,u_x$, $V=0$, zero terminal value, and the source
\[
 \frac{d\varphi'}2\Bigl[(v-v')_{xx}+a\bigl(v_x^2-v_x'^2\bigr)\Bigr]
 +\frac{c_\theta a}2\Bigl[(v_x-u_x)(D_\delta)_x-(v'_x-u_x)(D_{\delta'})_x\Bigr],
\]
where $v'=u_{\theta+\delta'}$. By the bounds just proved, every spatial
derivative of this source is bounded by $C_j(|\delta|+|\delta'|)$. By
the same induction, we have
$\sup_{q,x}|\partial_x^j(D_\delta-D_{\delta'})|\le C_j(|\delta|+|\delta'|)$.
It follows that, as $\delta\to0$, the functions $D_\delta$ converge
uniformly, together with all their spatial derivatives, to a function
$\dot u_\theta$. Passing to the limit in the integrated form of
\eqref{eq:sc-Dh}, we find that $\dot u_\theta$ solves
\begin{equation}\label{eq:sc-W}
 \partial_q\dot u_\theta+\frac{c_\theta}2(\dot u_\theta)_{xx}
   +c_\theta a\,u_{\theta,x}(\dot u_\theta)_x
 =-\frac{d\varphi'}2\bigl(u_{\theta,xx}+a\,u_{\theta,x}^2\bigr),
 \qquad \dot u_\theta(1,\cdot)=0 .
\end{equation}
We conclude that $u_\theta$ and each of its spatial derivatives are
differentiable in $\theta$, with
$\partial_\theta\partial_x^ju_\theta=\partial_x^j\dot u_\theta$, and that
these derivatives are bounded uniformly in $(q,x)$ and $\theta$.
To prove that these derivatives depend continuously on $\theta$, we
set $b_\theta=c_\theta a\,u_{\theta,x}$,
$F_\theta=(d\varphi'/2)(u_{\theta,xx}+a\,u_{\theta,x}^2)$, and
$R=\dot u_\theta-\dot u_{\theta'}$. Subtracting \eqref{eq:sc-W} at the
two parameters, we obtain
\[
 \Bigl(\partial_q+\frac{c_\theta}2\partial_{xx}+b_\theta\partial_x\Bigr)R
 =-(F_\theta-F_{\theta'})
  -\frac{c_\theta-c_{\theta'}}2(\dot u_{\theta'})_{xx}
  -(b_\theta-b_{\theta'})(\dot u_{\theta'})_x,
 \qquad R(1,\cdot)=0 .
\]
Since $c_\theta-c_{\theta'}=(\theta-\theta')d\varphi'$, and since
$u_\theta-u_{\theta'}$ and all its spatial derivatives are
$O(|\theta-\theta'|)$ by the bound on $\partial_\theta\partial_x^ju_\theta$,
every fixed-order spatial derivative of the right side is
$O(|\theta-\theta'|)$. Using the same backward estimates, we have
$\sup_{q,x}|\partial_x^jR|\le C_j|\theta-\theta'|$, so
$\partial_x^j\dot u_\theta$ depends continuously on $\theta$, uniformly
in $(q,x)$. The same argument applies without
change to transports $q\mapsto q+\theta_1\varphi_1(q)+\theta_2\varphi_2(q)$
with several parameters, and gives continuous partial derivatives in
each parameter.

\paragraph{Step 2: the variation of the solution.}
We write $\dot u=\dot u_0$ for the derivative at $\theta=0$. By
\eqref{eq:sc-W} and \eqref{eq:sc-backward}, with $Z=Y$, we have
\begin{equation}\label{eq:sc-udot-rep}
 \dot u(q,y)=\frac d2\int_q^1\varphi'(r)\Psi_q(r;y)\dd r,\qquad
 \Psi_q(r;y)=\E_{q,y}\bigl[h_r+a(r)m_r^2\bigr],
\end{equation}
where $\E_{q,y}$ refers to the diffusion started from $y$ at level $q$.
To evaluate $\Psi_q(r;y)$, we use Itô's formula along $Y$. Differentiating
\eqref{eq:sc-pde-integrated} twice in $x$ and using Itô's formula for
$u_{xx}(r,Y_r)$, we find
\[
 \dd h_r=-d\,a(r)h_r^2\dd r+\sqrt d\,u_{xxx}(r,Y_r)\dd B_r,
\]
and by \eqref{eq:sc-martingale},
$\dd(m_r^2)=d\,h_r^2\dd r+2\sqrt d\,m_rh_r\dd B_r$. The product rule for
the continuous semimartingale $m_r^2$ and the deterministic function
$a$ of bounded variation gives
$\dd(a(r)m_r^2)=a(r-)\dd(m_r^2)+m_r^2\,\mu(\dd r)$ on $(0,1]$, and
$a(r-)=a(r)$ except at countably many $r$. Adding the two
differentials, we see that the drift terms $\mp d\,a(r)h_r^2\dd r$
cancel, and
\[
 \dd\bigl(h_r+a(r)m_r^2\bigr)
 =\sqrt d\bigl[u_{xxx}(r,Y_r)+2a(r)m_rh_r\bigr]\dd B_r+m_r^2\,\mu(\dd r).
\]
Taking conditional expectations given $Y_q=y$, and recalling that
$\mu=\nu$ on $(q,r]\subset(0,1]$, we obtain
\begin{equation}\label{eq:sc-Fq}
 \Psi_q(r;y)-\Psi_q(q;y)=\int_{(q,r]}g_q^v(y)\,\nu(\dd v),\qquad q\le r\le1,
\end{equation}
where $\Psi_q(q;y)=u_{xx}(q,y)+a(q)u_x(q,y)^2$.

We insert \eqref{eq:sc-Fq} into \eqref{eq:sc-udot-rep} and exchange the
order of integration over $q<v\le r\le1$. Note that
$\int_q^1\varphi'(r)\dd r=-\varphi(q)$ and
$\int_v^1\varphi'(r)\dd r=-\varphi(v)$, since $\varphi(1)=0$. We obtain
\begin{equation}\label{eq:scalarNodeVariation}
 \dot u(q,y)=-\frac d2\Bigl[
  \varphi(q)\bigl(u_{xx}(q,y)+a(q)u_x(q,y)^2\bigr)
  +\int_{(q,1]}\varphi(r)g_q^r(y)\,\nu(\dd r)\Bigr].
\end{equation}
An atom at $q$ belongs to $a(q)$ and not to the integral over $(q,1]$.
An atom at $1$ belongs to the integral, but its factor $\varphi(1)$
vanishes.

The same computation holds at every base point $\theta$, with $Y^\theta$
in place of $Y$. The diffusion coefficient $d$ is then replaced by
$c_\theta(r)$; the drift terms $\mp c_\theta(r)a(r)h_r^2\dd r$ still
cancel, and the source in \eqref{eq:sc-W} still carries the factor
$d\varphi'$, which is the derivative of $c_\theta$ in $\theta$. At
$(q,y)=(0,0)$, we have $\varphi(0)=0$ and
$\E[\partial_xu_\theta(r,Y^\theta_r)^2\mid Y^\theta_0=0]
=\widetilde\Gamma_\theta(r)$. This gives
\begin{equation}\label{eq:sc-root-variation}
 \frac{\dd}{\dd\theta}u_\theta(0,0)
 =-\frac d2\int\varphi(r)\widetilde\Gamma_\theta(r)\,\nu(\dd r),
 \qquad|\theta|<\theta_0 .
\end{equation}

\paragraph{Step 3: Gaussian differentiation at the fixed profile.}
Let $B$ be a standard Brownian motion, and under the reference law let
$X^\theta_q=B_{dT_\theta(q)}$. This Gaussian path has covariance
$K_\theta(q,r)=d\,T_\theta(q\wedge r)$ and the law of
$\int_0^q\sqrt{c_\theta(r)}\dd B_r$. Then \cref{lem:sc-likelihood},
applied with the coefficient $c_\theta$, implies that the joint law of $(Y^\theta,\epsilon)$
has density \eqref{eq:sc-joint-law} with $u_\theta$ in place of $u$,
relative to the law of $X^\theta$ and counting measure. The factor
$e^{-su_\theta(0,0)}$ does not depend on the path and cancels from
normalized expectations. With
\[
 \ell_\theta(X,\epsilon)=\epsilon X_1-\int u_\theta(r,X_r)\,\nu(\dd r),
 \qquad\Phi_\theta(X)=u_{\theta,x}(p,X_p)^2,
\]
we have
$\widetilde\Gamma_\theta(p)=\E_\theta\Phi_\theta$, where $\E_\theta$ is
the normalized expectation with weight $W_\theta=e^{\ell_\theta}$ in
the notation of \cref{lem:fnd-path-heat}. At $\theta=0$, the path $X$
under $\E_0$ has the law of $Y$. As in \cref{sec:sc-law}, we write
$m_r=u_x(r,X_r)$ and $h_r=u_{xx}(r,X_r)$, and we also write
$u_{3,r}=u_{xxx}(r,X_r)$.

We check that the assumptions (A1)--(A4) of \cref{lem:fnd-path-heat}
hold, with $z=\theta$ in the compact interval
$I=[-\theta_0/2,\theta_0/2]$. The Gaussian coordinates consist of one
field path on one edge, with the deterministic clock $dT_\theta(q)$,
which is continuous in $(q,\theta)$, continuously differentiable in
$\theta$, and strictly increasing in $q$; there are no finite
coordinates, and the terminal spin $\epsilon$ is the counting variable.
This is (A1). The kernels $K_\theta$ and
$\dot K_\theta(q,r)=d\varphi(q\wedge r)$ are jointly continuous in
$(\theta,q,r)$, which is (A2). In the notation of
\cref{lem:fnd-path-heat}, where $M$ and $Q$ denote derivative measures,
the derivative measures of $\ell_\theta$ and $\Phi_\theta$ are
\[
 \begin{aligned}
 S&=\epsilon\delta_1-u_{\theta,x}(r,X_r)\,\nu(\dd r),&\qquad
 M&=-\iota_\#\bigl(u_{\theta,xx}(r,X_r)\,\nu(\dd r)\bigr),\\
 P&=2u_{\theta,x}u_{\theta,xx}(p,X_p)\,\delta_p,&\qquad
 Q&=2\bigl(u_{\theta,xx}^2+u_{\theta,x}u_{\theta,xxx}\bigr)(p,X_p)\,
   \delta_{(p,p)},
 \end{aligned}
\]
where $\iota(r)=(r,r)$. The measure $M$ is carried by the diagonal,
with an atom at $(r,r)$ for each atom $r$ of $\nu$, as in
\cref{rem:gauss-local}. Using \cref{lem:sc-basic}(b), we see that these
measures have
total variation bounded uniformly in $\theta$ and $X$, and by Step 1,
they are continuous in $(\theta,X)$ in total variation, where $X$
carries the supremum norm. The maps $X\mapsto\ell_\theta(X,\epsilon)$
and $X\mapsto\Phi_\theta(X)$ are twice continuously differentiable,
since $u_\theta$ has bounded continuous spatial derivatives. This is
(A4). The observable $\Phi_\theta$ takes values in $[0,1]$, and the
explicit parameter derivatives
\[
 \dot\Phi_\theta=2u_{\theta,x}(p,X_p)\,\partial_x\dot u_\theta(p,X_p),\qquad
 \dot\ell_\theta=-\int\dot u_\theta(r,X_r)\,\nu(\dd r)
\]
are bounded and jointly continuous in $(\theta,X)$ by Step 1. Since
$|u_{\theta,x}|\le1$ and $u_\theta(r,0)$ is bounded, we have
$|\ell_\theta|\le C+2\sup_r|X_r|$, which is the linear growth required
in (A3). We conclude that \cref{lem:fnd-path-heat} applies to
$\E_\theta\Phi_\theta$. By \cref{rem:gauss-local}, the right side of
\eqref{eq:fnd-path-heat} is continuous in $\theta$, which proves that
$\theta\mapsto\widetilde\Gamma_\theta(p)$ is continuously differentiable
on $(-\theta_0/2,\theta_0/2)$. Each term on the right side is bounded
by a constant depending only on $\beta$ and $\varphi$; in particular,
the derivative is bounded uniformly in $p\in[0,1]$ and $\theta$. We use
this uniform bound in the proof of \cref{lem:Jdom}.

\paragraph{Step 4: the terms of the identity.}
We evaluate the right side of \eqref{eq:fnd-path-heat} at $\theta=0$
and divide it by $d$. For a signed measure $R$ on $[0,1]^2$, we have
$d^{-1}\langle\dot K_0,R\rangle=\iint\varphi(q\wedge r)\,R(\dd q,\dd r)$.
Note that, under $\E_0$, the centered observable is
$\Phi_0-\E_0\Phi_0=\bar f_p$.
We group the contributions as follows.
\begin{enumerate}
\item The second derivative of the observable contributes
 \[
  \mathrm I_1=\frac1{2d}\E_0\langle\dot K_0,Q\rangle
  =\varphi(p)\E^0\bigl[h_p^2+m_pu_{3,p}\bigr].
 \]
\item The explicit derivative of the observable is
 $\dot\Phi_0=2m_p\partial_y\dot u(p,X_p)$. Differentiating
 \eqref{eq:scalarNodeVariation} in $y$ under the integral, which is
 allowed by \eqref{eq:sc-dg}, we find
 \[
  \mathrm I_2=\frac1d\E_0\dot\Phi_0
  =-\varphi(p)\E^0\bigl[m_pu_{3,p}+2a(p)m_p^2h_p\bigr]
  -\int_{(p,1]}\varphi(\ell)\E^0\bigl[m_p\partial_yg_p^\ell(Y_p)\bigr]
   \,\nu(\dd\ell).
 \]
\item Since $\dot K_0$ is symmetric and $\dot K_0(1,p)=d\varphi(p)$,
 the score--observable terms give
 \[
  \begin{aligned}
  \mathrm I_3&=\frac1{2d}\E_0\langle\dot K_0,S\otimes P+P\otimes S\rangle
  =\frac1d\E_0\langle\dot K_0,S\otimes P\rangle\\
  &=2\varphi(p)\E^0[m_ph_p\epsilon]
  -2\int\varphi(p\wedge r)\E^0[m_ph_pm_r]\,\nu(\dd r).
  \end{aligned}
 \]
\item The product of the scores contributes the following; the
 terminal--terminal term vanishes, since $\varphi(1)=0$:
 \[
  \begin{aligned}
  \mathrm I_4&=\frac1{2d}\E_0\bigl[\bar f_p\langle\dot K_0,S\otimes S\rangle\bigr]\\
  &=-\int\varphi(r)\E^0[\bar f_p\epsilon m_r]\,\nu(\dd r)
  +\frac12\iint\varphi(r\wedge\ell)\E^0[\bar f_pm_rm_\ell]
   \,\nu(\dd r)\nu(\dd\ell).
  \end{aligned}
 \]
\item The Hessian of the logarithmic weight contributes
 \[
  \mathrm I_5=\frac1{2d}\E_0\bigl[\bar f_p\langle\dot K_0,M\rangle\bigr]
  =-\frac12\int\varphi(r)\E^0[\bar f_ph_r]\,\nu(\dd r).
 \]
\item The explicit derivative of the weight is
 $\dot\ell_0=-\int\dot u(r,X_r)\,\nu(\dd r)$. Evaluating it with
 \eqref{eq:scalarNodeVariation}, we find
 \[
  \begin{aligned}
  \mathrm I_6=\frac1d\Cov_0(\Phi_0,\dot\ell_0)
  ={}&\frac12\int\varphi(r)\E^0\bigl[\bar f_p\bigl(h_r+a(r)m_r^2\bigr)\bigr]
   \,\nu(\dd r)\\
  &+\frac12\int\nu(\dd r)\int_{(r,1]}\varphi(\ell)
   \E^0\bigl[\bar f_pg_r^\ell(Y_r)\bigr]\,\nu(\dd\ell).
  \end{aligned}
 \]
\end{enumerate}
By \eqref{eq:fnd-path-heat}, we have
\[
 \frac1d\,\frac{\dd}{\dd\theta}\widetilde\Gamma_\theta(p)\Big|_{\theta=0}
 =\mathrm I_1+\cdots+\mathrm I_6 .
\]
Conditioning on $\calG^0_p$ and using \eqref{eq:sc-martingale}, we
have
$\E^0[m_ph_p\epsilon]=\E^0[m_p^2h_p]$. The following table records
where each part of $\mathrm I_1,\ldots,\mathrm I_6$ goes; ``line $i$''
refers to the $i$-th line of \eqref{eq:scalarJexplicit}.
\begin{center}
\small
\begin{tabular}{@{}lll@{}}
\toprule
Group & Part & Contribution\\
\midrule
$\mathrm I_1$ & $\varphi(p)\E^0h_p^2$ & line 1, diagonal coefficient\\
 & $\varphi(p)\E^0[m_pu_{3,p}]$ & cancels with $\mathrm I_2$\\
$\mathrm I_2$ & $-\varphi(p)\E^0[m_pu_{3,p}]$ & cancels with $\mathrm I_1$\\
 & $-2a(p)\varphi(p)\E^0[m_p^2h_p]$ & line 1, diagonal coefficient\\
 & integral over $(p,1]$ & line 1, second term\\
$\mathrm I_3$ & $2\varphi(p)\E^0[m_ph_p\epsilon]=2\varphi(p)\E^0[m_p^2h_p]$
 & line 1, diagonal coefficient\\
 & integral of $\E^0[m_ph_pm_r]$ & line 2, first term\\
$\mathrm I_4$ & integral of $\E^0[\bar f_p\epsilon m_r]$ & line 2, second term\\
 & double integral & line 3, second term\\
$\mathrm I_5$ & $-(1/2)\int\varphi(r)\E^0[\bar f_ph_r]\,\nu(\dd r)$
 & cancels with $\mathrm I_6$\\
$\mathrm I_6$ & $(1/2)\int\varphi(r)\E^0[\bar f_ph_r]\,\nu(\dd r)$
 & cancels with $\mathrm I_5$\\
 & $(1/2)\int\varphi(r)a(r)\E^0[\bar f_pm_r^2]\,\nu(\dd r)$
 & line 3, first term\\
 & iterated integral & line 4\\
\bottomrule
\end{tabular}
\end{center}
In summary, the terms with $m_pu_{3,p}$ cancel, the covariance terms with
$h_r$ cancel, and the diagonal coefficient is
$\E^0[h_p^2-2a(p)m_p^2h_p+2m_p^2h_p]=\E^0[h_p^2+2(1-a(p))m_p^2h_p]$.
The sum $\mathrm I_1+\cdots+\mathrm I_6$ is the right side of
\eqref{eq:scalarJexplicit}, including its signs, its factors $1/2$, and
its strict future intervals. This completes the proof.
\end{proof}

\section{Domination by the scalar multiplier}\label{sec:sc-domination}

We now show that the scalar response operator is dominated by its
multiplier.
We use the scalar Parisi functional with coefficient $d$,
\[
 \mathcal P_d(a)=u_{a,d}(0,0)-\frac d2\int_0^1q\,a(q)\dd q,
\]
defined on the convex set of profiles. For $d=\beta^2$, it is the
functional $\mathcal P_\beta$ of \eqref{eq:intro-parisi}.

\begin{lemma}[Scalar multiplier domination]\label{lem:Jdom}
For every profile $a$ and $d\in[\beta^2/2,\beta^2]$, the operator
$J^{\scal}_{a,d}$ is self-adjoint on $\mathcal H_\nu$, with norm at most
$C_\beta$, and
\begin{equation}\label{eq:Jdom}
 J^{\scal}_{a,d}\preceq M_{k_{a,d}},
 \qquad (M_{k_{a,d}}v)(q)=k_{a,d}(q)v(q).
\end{equation}
\end{lemma}

\begin{proof}
The bound on the norm is \eqref{eq:sc-J-bounded}. The proof of the rest
has five steps. We first recall that $\mathcal P_d$ is convex and
compute its first variation. Then we follow $\mathcal P_d$ along a
transport of levels and compute its second derivative with
\cref{lem:sc-response}. Next, we bound the same second derivative from
below, using convexity with the first variation frozen at $a$; the
comparison yields the inequality \eqref{eq:Jdom} for smooth directions.
Two-parameter transports then give the symmetry, and a density argument
extends both conclusions to $\mathcal H_\nu$. No minimizing property of
$a$ is used. We abbreviate $\Gamma=\Gamma_{a,d}$ and $k=k_{a,d}$, and
$\langle\cdot,\cdot\rangle$ denotes the inner product of $L^2(\nu)$.

\paragraph{Step 1: convexity and the first variation.}
We first note that $\mathcal P_d$ is convex on profiles. Indeed, by
\cite[Theorem~2]{AuffingerChen2015Uniqueness}, applied with
$\xi(q)=dq^2/2$ and zero external field, the value at $(0,0)$ of the
solution with terminal value $\log\cosh x$ is a strictly convex
function of the measure, and hence of its distribution function. This
value differs from $u_{a,d}(0,0)$ by $\log2$, and the penalty term of
$\mathcal P_d$ is affine in $a$.

The first variation of $\mathcal P_d$ along segments is the scalar case of
\cite[(2.8)]{AL} (see also \cite[Lemma~3.2]{JagannathTobasco2017PhaseDiagram}
and \cite[Theorem~2]{Chen2017Variational}): for profiles $a$ and
$\widetilde a$,
\begin{equation}\label{eq:sc-first-variation}
 \lim_{\vartheta\downarrow0}\frac{\mathcal P_d(a+\vartheta(\widetilde a-a))
                               -\mathcal P_d(a)}{\vartheta}
 =\frac d2\int_0^1(\Gamma(q)-q)(\widetilde a(q)-a(q))\dd q .
\end{equation}
The statement in \cite{AL} assumes a positive root mass. We note
that the argument of Step 1 of the proof of \cref{lem:sc-response},
with the coefficient fixed and the profile varied, gives a direct proof
of \eqref{eq:sc-first-variation} for all profiles. The difference
quotients of
$u_{a+\vartheta(\widetilde a-a),d}$ converge, with all spatial
derivatives, to the solution $\dot u$ of
\[
 \partial_q\dot u+\frac d2\dot u_{xx}+d\,a\,u_x\dot u_x
 =-\frac d2(\widetilde a-a)u_x^2,\qquad \dot u(1,\cdot)=0,
\]
and \eqref{eq:sc-backward} gives
\[
 \dot u(0,0)=\frac d2\int_0^1(\widetilde a(q)-a(q))\,\E^0m_q^2\dd q
 =\frac d2\int_0^1(\widetilde a(q)-a(q))\Gamma(q)\dd q .
\]
The penalty contributes $-(d/2)\int_0^1q(\widetilde a(q)-a(q))\dd q$.
Since $\mathcal P_d$ is convex, the difference quotient in
\eqref{eq:sc-first-variation} is nondecreasing in $\vartheta\in(0,1]$.
Comparing its value at $\vartheta=1$ with its limit as
$\vartheta\downarrow0$, we obtain
\begin{equation}\label{eq:sc-convexity}
 \mathcal P_d(\widetilde a)-\mathcal P_d(a)
 \ge\frac d2\int_0^1(\Gamma(q)-q)(\widetilde a(q)-a(q))\dd q .
\end{equation}

\paragraph{Step 2: the second derivative along a transport.}
Let $\varphi$, $T_\theta$, $\theta_0$, and $a_\theta=a\circ T_\theta^{-1}$
be as in \cref{sec:sc-response}, and let
$\Pi(\theta)=\mathcal P_d(a_\theta)$ for $|\theta|<\theta_0/2$. By
\eqref{eq:sc-transported} at $q=0$, we have
$u_{a_\theta,d}(0,0)=u_\theta(0,0)$. Using Fubini's theorem, we have
\[
 \int_0^1x\,a_\theta(x)\dd x
 =\int\Bigl(\int_{T_\theta(q)}^1x\dd x\Bigr)\mu(\dd q)
 =\frac12-\frac12\int T_\theta(q)^2\,\mu(\dd q),
\]
which gives
\[
 \Pi(\theta)=u_\theta(0,0)-\frac d4+\frac d4\int T_\theta^2\dd\mu .
\]
The derivative of the last term is $(d/2)\int T_\theta\varphi\dd\mu$,
in which the atom $s\delta_0$ does not contribute, since $\varphi(0)=0$.
Combining this with \eqref{eq:sc-root-variation}, we obtain
\begin{equation}\label{eq:transporthess}
 \Pi'(\theta)=\frac d2\int\varphi\bigl(T_\theta-\widetilde\Gamma_\theta\bigr)\dd\nu,
 \qquad
 \Pi''(0)=\frac d2\langle\varphi,(I-J_{a,d})\varphi\rangle .
\end{equation}
For the second identity, we used \cref{lem:sc-response} and
differentiated under the integral over $\nu$. This is justified by
dominated convergence, since by Step 3 of the proof of
\cref{lem:sc-response}, $\partial_\theta\widetilde\Gamma_\theta(p)$ is
continuous in $\theta$ and bounded uniformly in $p$ and $\theta$. By
the same argument, $\Pi''$ is continuous near $0$. At $\theta=0$, the first identity
in \eqref{eq:transporthess} is the first variation along transports in
\cite[Section~3, (19)]{AuffingerChen2015Properties}, where it is
credited to \cite[Lemma~3.7, Eq.~(3.34)]{Talagrand2006ParisiMeasures}.

\paragraph{Step 3: the lower bound from convexity.}
Let $G(y)=\int_y^1(\Gamma(q)-q)\dd q$. Recall that
$a_\theta(q)=\mu(\{y:T_\theta(y)\le q\})$. By Fubini's theorem, we have
$\int_0^1(\Gamma(q)-q)a_\theta(q)\dd q=\int G(T_\theta y)\,\mu(\dd y)$.
Applying \eqref{eq:sc-convexity} with $\widetilde a=a_\theta$, we obtain
\[
 \Pi(\theta)-\Pi(0)\ge\frac d2\int\bigl(G(T_\theta y)-G(y)\bigr)\mu(\dd y)
 =\frac d2\int\bigl(G(T_\theta y)-G(y)\bigr)\nu(\dd y),
\]
where the atom $s\delta_0$ does not contribute because $T_\theta(0)=0$.
Then \cref{lem:sc-basic}(c) implies that $G$ is twice continuously
differentiable on $[0,1]$, with
\[
 G'(y)=y-\Gamma(y),\qquad G''(y)=1-k(y).
\]
The function $G''$ is uniformly continuous and $\varphi$ is bounded, so
by Taylor's formula, we have
\[
 G(T_\theta y)-G(y)=\theta\varphi(y)G'(y)
 +\frac{\theta^2}2\varphi(y)^2\bigl(1-k(y)\bigr)+o(\theta^2),
\]
uniformly in $y\in[0,1]$. By \eqref{eq:transporthess} at $\theta=0$,
we have $\Pi'(0)=(d/2)\int\varphi G'\dd\nu$. Integrating the expansion
against $\nu$ and inserting it into the previous inequality, we obtain
\[
 \Pi(\theta)-\Pi(0)-\theta\Pi'(0)
 \ge\frac{\theta^2}2\cdot\frac d2\int(1-k)\varphi^2\dd\nu+o(\theta^2).
\]
On the other hand, $\Pi$ is twice continuously differentiable near
$0$, so the left side equals
$(\theta^2/2)\Pi''(0)+o(\theta^2)$. Dividing by $\theta^2/2$ and
letting $\theta\to0$, we find
\[
 \Pi''(0)\ge\frac d2\int(1-k)\varphi^2\dd\nu .
\]
Comparing this with
\eqref{eq:transporthess}, we conclude that
\begin{equation}\label{eq:sc-dom-smooth}
 \langle\varphi,J_{a,d}\varphi\rangle\le\langle\varphi,k\varphi\rangle .
\end{equation}
Note that $\Pi$ need not be convex, since $\theta\mapsto a_\theta$ is
not affine. The lower bound accounts for this through the term
$1-k=G''$.

\paragraph{Step 4: symmetry.}
We now take two functions $\varphi_1$ and $\varphi_2$ as above, and we
consider the transport $T(q)=q+\theta_1\varphi_1(q)+\theta_2\varphi_2(q)$, which
fixes both endpoints. Let $\widetilde\Gamma_{\theta_1,\theta_2}$ and
$\Pi(\theta_1,\theta_2)$ be defined as before, for small
$|\theta_1|,|\theta_2|$. By Step 1 of the proof of
\cref{lem:sc-response} for several parameters, and by
\cref{lem:fnd-path-heat} applied in each parameter with the other one
fixed, the function $\widetilde\Gamma_{\theta_1,\theta_2}(p)$ has
continuous first partial derivatives, bounded uniformly in $p$. The
computation of Step 2 of that proof at every base point, with the
coefficient derivative $d\varphi_i'$, yields
\[
 \partial_{\theta_i}\Pi
 =\frac d2\int\varphi_i\bigl(T-\widetilde\Gamma_{\theta_1,\theta_2}\bigr)\dd\nu,
 \qquad i=1,2 .
\]
It follows that $\Pi$ has continuous second partial derivatives, and
using \cref{lem:sc-response} in the direction $\varphi_2$ (with
$\theta_1=0$), we find that
\[
 \partial_{\theta_2}\partial_{\theta_1}\Pi(0,0)
 =\frac d2\langle\varphi_1,(I-J_{a,d})\varphi_2\rangle .
\]
By the equality of mixed partial derivatives, we conclude that this
bilinear form is symmetric, that is, $\langle\varphi_1,J_{a,d}\varphi_2\rangle
=\langle\varphi_2,J_{a,d}\varphi_1\rangle$.

\paragraph{Step 5: extension to the transport space.}
In Steps 3 and 4, we proved the domination inequality and the symmetry
for restrictions to $[0,1]$ of smooth functions vanishing at both
endpoints. Such functions lie in $\mathcal H_\nu$, and they are dense
in it. Indeed, let $v\in\mathcal H_\nu$, and let $\psi_n$ be
continuous functions with $\|\psi_n-v\|\to0$, which exist because
continuous functions are dense in $L^2(\nu)$. For $n\ge3$, let $\chi_n$ be
continuous with values in $[0,1]$, vanishing at $0$ and $1$, and equal
to one on $[1/n,1-1/n]$. Then
$\|\chi_n\psi_n-v\|\le\|\psi_n-v\|+\|(1-\chi_n)v\|$, and
$\|(1-\chi_n)v\|\to\|\ind_{\{0,1\}}v\|=0$ by dominated convergence,
because $v$ vanishes $\nu$-almost everywhere on $\{0,1\}$. Finally, let
$f$ be continuous with $f(0)=f(1)=0$. If $P$ is a polynomial with
$|P-f|\le\tau$ on $[0,1]$, then the polynomial
$P(x)-P(0)(1-x)-P(1)x$ vanishes at both endpoints and is within
$3\tau$ of $f$; by the Weierstrass theorem, $f$ is a uniform limit of
such polynomials. For $v\in\mathcal H_\nu$, we have
$\langle v,J^{\scal}_{a,d}v\rangle=\langle v,J_{a,d}v\rangle$, and both
sides of \eqref{eq:sc-dom-smooth} and of the symmetry identity are
continuous on $L^2(\nu)$, since $J_{a,d}$ is bounded on $L^2(\nu)$ (see
the proof of \eqref{eq:sc-J-bounded}) and $0\le k\le\beta^2$. Both conclusions extend to $\mathcal H_\nu$ by
continuity. Multiplication by $k$ maps $\mathcal H_\nu$ into
itself, and \eqref{eq:Jdom} follows.
\end{proof}

\section{The reference multiplier}\label{sec:sc-reference}

A finite-dimensional minimizer may place mass outside the scalar
support. For this reason, we control the multiplier at the scalar
Parisi measure on all of $[0,1]$. Recall from
\eqref{eq:loc-scalar-measure} that $\supp\mu_\beta=[0,q_\beta]$ with
$0<q_\beta<1$ \cite[Theorem~1.1]{Lopatto}, and that $\mu_\beta$ has an
atom at $q_\beta$. In particular, $\alpha_\beta=1$ on $[q_\beta,1]$.

\begin{lemma}[Reference multiplier on the whole interval]\label{lem:reference}
At the scalar Parisi measure,
\begin{equation}\label{eq:referencek}
 k_{\alpha_\beta,\beta^2}(q)=1\quad(0\le q\le q_\beta),
 \qquad
 k_{\alpha_\beta,\beta^2}(q)
       \le e^{-\beta^2(q-q_\beta)/2}\quad(q_\beta\le q\le1).
\end{equation}
In particular, $\sup_{q\in[0,1]}k_{\alpha_\beta,\beta^2}(q)\le1$.
\end{lemma}

\begin{proof}
We abbreviate $d=\beta^2$, $a=\alpha_\beta$, $k=k_{\alpha_\beta,d}$,
and $\Gamma=\Gamma_{\alpha_\beta,d}$, and we write $u$ and $Y$ for the
corresponding solution and diffusion. The proof has four steps. We
first prove that $k=1$ on the support. We then recall that a
transformed density of $Y_{q_\beta}$ is even and log-concave, show that
this density evolves by the heat equation above $q_\beta$, and finally
use a rearrangement inequality to bound $k$ there.

\paragraph{Step 1: the support.}
By an identity of Talagrand \cite{Talagrand2006ParisiMeasures}, in the
form stated in \cite[Theorem~5]{AuffingerChen2015Properties} and
\cite[Lemma~2.6]{Lopatto} (see also
\cite[Proposition~1.1]{JagannathTobasco2017PhaseDiagram} and
\cite[Proposition~1]{Chen2017Variational}), we have $\Gamma(q)=q$ for
every $q\in\supp\mu_\beta=[0,q_\beta]$. The function $\Gamma$ of
\cite{Lopatto} coincides with ours, since $u_x$ does not depend on the constant
$\log2$ in the terminal value. By \cref{lem:sc-basic}(c), we have
$\Gamma'=k$, which implies that $k=1$ on $(0,q_\beta)$. Since $k$ is
continuous, we conclude that $k=1$ on $[0,q_\beta]$.

\paragraph{Step 2: log-concavity of the transformed density.}
For $q\in(0,1]$, let $\pi_q$ be the density of $Y_q$, and set
\[
 z_q(x)=\pi_q(x)e^{-a(q)u(q,x)},\qquad
 V_q(x)=-\log z_q(x)=a(q)u(q,x)-\log\pi_q(x).
\]
By \eqref{eq:sc-prefix-law}, $z_q(x)\dd x$ is the distribution of the endpoint
$X_q$ under the finite measure obtained by weighting the law of the
Gaussian path, without normalization, by
$\exp\{-\int_{[0,q]}u(r,X_r)\,\mu_\beta(\dd r)\}$. By
\cite[Proposition~4.1]{Lopatto}, which applies to every probability
measure on $[0,1]$, including measures with atoms, the density $\pi_q$
is smooth and positive, and the potential $a(q)u^{\rm L}(q,x)-\log\pi_q(x)$
is smooth, even, and has second derivative at least $1/(dq)$; here
$u^{\rm L}=u-\log2$ is the solution with the terminal value of
\cite{Lopatto}, and $a(q)$ is the right-continuous value. The two
potentials differ by the constant $a(q)\log2$, which implies that
\[
 V_q''(x)\ge\frac1{dq}>0,\qquad V_q(-x)=V_q(x),\qquad x\in\R .
\]
Note that $V_q'(0)=0$ by evenness. Then $V_q'(x)\ge x/(dq)>0$ for
$x>0$.
This implies that $z_q$ is even, log-concave, and decreasing on
$(0,\infty)$. We use this at $q=q_\beta$, where the right-continuous
value is $a(q_\beta)=1$, including the atom of $\mu_\beta$ at
$q_\beta$; in particular, $z_{q_\beta}=\pi_{q_\beta}e^{-u(q_\beta,\cdot)}$.

\paragraph{Step 3: heat evolution above the support.}
On $[q_\beta,1]$, we have $a=1$, and the solution of \eqref{eq:sc-pde}
is
\[
 u(q,x)=\log(2\cosh x)+\frac{d(1-q)}2,\qquad q_\beta\le q\le1,
\]
by \cite[Lemma~2.4(i)]{Lopatto}; this formula can also be checked by
direct substitution, since $u_x=\tanh x$, $u_{xx}=\sech^2x$,
$u_{xx}+u_x^2=1$, and $\partial_qu=-d/2$. We write $G_v$ for the centered
Gaussian density with variance $v$. By \cite[Lemma~2.4(ii)]{Lopatto},
applied with mass $m=1$ on $[q_\beta,1)$, we have, for
$q_\beta<q\le1$,
\[
 \pi_q(x)=e^{u(q,x)}\int G_{d(q-q_\beta)}(x-y)\,e^{-u(q_\beta,y)}\pi_{q_\beta}(y)\dd y ;
\]
the constants $\log2$ in $u(q,x)$ and $u(q_\beta,y)$ cancel. This means
that
\begin{equation}\label{eq:sc-heat}
 z_q=G_{d(q-q_\beta)}*z_{q_\beta},\qquad q_\beta<q\le1 .
\end{equation}
Equivalently, $z_q$ solves the heat equation
$\partial_qz_q=(d/2)\partial_{xx}z_q$. To see this directly, we
substitute $\pi_q=e^{u}z_q$ into the forward equation of $Y$ on the
plateau, which gives
\[
 \partial_qz_q+(\partial_qu)z_q
 =e^{-u}\Bigl(\frac d2\partial_{xx}\pi_q-d\,\partial_x(u_x\pi_q)\Bigr)
 =\frac d2\partial_{xx}z_q-\frac d2\bigl(u_{xx}+u_x^2\bigr)z_q ;
\]
the zeroth-order terms cancel because $u_{xx}+u_x^2=1$ and
$\partial_qu=-d/2$.

On the plateau, we have $u_{xx}=\sech^2$ and
$\pi_q=e^{u(q,\cdot)}z_q=2\cosh(\cdot)\,e^{d(1-q)/2}z_q$. The multiplier
is therefore
\begin{equation}\label{eq:sc-k-plateau}
 k(q)=d\int\pi_q(x)\sech^4x\dd x
 =2d\,e^{d(1-q)/2}\int_\R z_q(x)\sech^3x\dd x,\qquad q_\beta\le q\le1 .
\end{equation}

\paragraph{Step 4: the rearrangement inequality.}
Let $f$ and $g$ be nonnegative, even, continuous functions that are
nonincreasing on $[0,\infty)$. We claim that
\begin{equation}\label{eq:sc-rearrangement}
 \int f(x-y)g(x)\dd x\le\int f(x)g(x)\dd x,\qquad y\in\R .
\end{equation}
This is the centered rearrangement inequality in one dimension, which
also follows from Anderson's inequality~\cite[Theorem~1]{Anderson1955}
by the layer-cake representation. We include a proof based on intersections of intervals.
By the layer-cake representation and Tonelli's theorem,
\[
 \int f(x-y)g(x)\dd x=\int_0^\infty\!\!\int_0^\infty
 \bigl|(\{f>t\}+y)\cap\{g>t'\}\bigr|\dd t\dd t',
\]
where $|\cdot|$ is Lebesgue measure. The superlevel sets $\{f>t\}$ and
$\{g>t'\}$ are open intervals centered at $0$. For two centered intervals
$I$ and $I'$, we have $|(I+y)\cap I'|\le\min\{|I|,|I'|\}=|I\cap I'|$, and
\eqref{eq:sc-rearrangement} follows.

We use \eqref{eq:sc-rearrangement} with $f=z_{q_\beta}$ and
$g=\sech^3$, which satisfy its hypotheses by Step 2. All the integrals
below are finite, since $\sech\le1$ and
$\int z_{q_\beta}=\E[e^{-u(q_\beta,Y_{q_\beta})}]\le1/2$. By
\eqref{eq:sc-heat} and Tonelli's theorem, for $q_\beta<q\le1$,
\[
 \int z_q(x)\sech^3x\dd x
 =\int G_{d(q-q_\beta)}(y)\Bigl(\int z_{q_\beta}(x-y)\sech^3x\dd x\Bigr)\dd y
 \le\int z_{q_\beta}(x)\sech^3x\dd x .
\]
Inserting this into \eqref{eq:sc-k-plateau}, and using
\eqref{eq:sc-k-plateau} at $q=q_\beta$ and $k(q_\beta)=1$ from Step 1,
we find
\[
 k(q)\le2d\,e^{d(1-q)/2}\int z_{q_\beta}(x)\sech^3x\dd x
 =e^{-d(q-q_\beta)/2}k(q_\beta)=e^{-d(q-q_\beta)/2},
 \qquad q_\beta\le q\le1 .
\]
Together with Step 1, this proves \eqref{eq:referencek}.
\end{proof}

\section{Stability at a minimizing profile}\label{sec:sc-minimizer}

The isolation estimate of \cref{prop:loc-isolation} compares a
minimizer with the scalar trial profile
$\gamma^{\rm tr}_\lambda=\max\{\lambda,\gamma_\beta\}$, where
$\gamma_\beta=\beta\alpha_\beta$. We first bound the distance between
the trial profile and $\gamma_\beta$. Recall from \eqref{eq:lz-gsroot}
that $\gamma_\beta(q)\ge k_0q$ for $0\le q\le T_0$, with constants
$k_0,T_0>0$ depending only on $\beta$. We claim that, for every
$\lambda>0$,
\begin{equation}\label{eq:sc-floor}
 \|\gamma^{\rm tr}_\lambda-\gamma_\beta\|_{L^1(0,1)}
 \le C_\beta\lambda^2 .
\end{equation}
Indeed, $\gamma^{\rm tr}_\lambda-\gamma_\beta=(\lambda-\gamma_\beta)_+$.
Suppose first that $\lambda\le k_0T_0$. For $q\ge T_0$, since
$\gamma_\beta$ is nondecreasing, we have
$\gamma_\beta(q)\ge\gamma_\beta(T_0)\ge k_0T_0\ge\lambda$, and the
difference vanishes there. This gives
\[
 \|\gamma^{\rm tr}_\lambda-\gamma_\beta\|_{L^1(0,1)}
 =\int_0^{T_0}(\lambda-\gamma_\beta(q))_+\dd q
 \le\int_0^{\lambda/k_0}(\lambda-k_0q)\dd q=\frac{\lambda^2}{2k_0}.
\]
If $\lambda>k_0T_0$, the difference is at most $\lambda$ everywhere,
and $\lambda\le\lambda^2/(k_0T_0)$. Since $T_0\le1/4<2$, both cases give
\eqref{eq:sc-floor} with $C_\beta=1/(k_0T_0)$.

\begin{proposition}[Scalar stability at the minimizing profile]
\label{prop:lz-stability}
There exists a constant $C_{\rm st}\ge1$, depending only on $\beta$ (in
particular, not on the minimizer or its atoms), with the following
property. Let $N\ge2$, $0<\rho\le1$, $t=1-\rho$, and
$0<\lambda\le\lambda_0$, and set $s=\lambda/\beta$. Let
$\mu_N=s\delta_0+\nu_N$ be a minimizer of $\mathcal F_{t,s}$ in
\eqref{eq:fnd-functional}, let $a_N(q)=\mu_N([0,q])$ be its profile,
and let $\eta$ be a trial budget for the minimizer $\beta a_N$ of
$F_\rho$, in the sense of \eqref{eq:loc-trialbudget}. Let
$d_N=\beta^2(1-t/N)$. Then
\begin{equation}\label{eq:lz-profileconvergence}
 \|a_N-\alpha_\beta\|_{L^1(0,1)}
 \le\min\bigl\{1,C_{\rm st}[\rho^{-3/4}\eta^{1/4}+\lambda^2]\bigr\}.
\end{equation}
Further, let
$\varepsilon_N=\max\{0,\sup_{q\in[0,1]}k_{a_N,d_N}(q)-1\}$. Then, on the
scalar transport space $\mathcal H_{\nu_N}$ of $\mu_N$,
\begin{equation}\label{eq:lz-stability}
 J^{\scal}_{a_N,d_N}\preceq(1+\varepsilon_N)I,
 \qquad
 \varepsilon_N\le C_{\rm st}
       \bigl[\rho^{-3/4}\eta^{1/4}+\lambda^2+N^{-1}\bigr].
\end{equation}
\end{proposition}

\begin{proof}
We use the energy normalization of \cref{sec:fnd}, in which
$\gamma=\beta a$ and $F_\rho(\gamma)=\beta^{-1}\mathcal F_{t,s}(\mu)$.
The map $\mu\mapsto\beta a$ is a bijection onto the admissible profiles
of \cref{sec:lz}, and multiplication by $\beta^{-1}$ preserves
minimizers. Hence $\gamma_N=\beta a_N$ is a minimizer of $F_\rho$, and
$\eta$ is a trial budget for it. The hypotheses of
\cref{prop:loc-isolation} hold, and \eqref{eq:loc-isolation} implies
that
$\|\gamma_N-\gamma^{\rm tr}_\lambda\|_{L^1}\le C\rho^{-3/4}\eta^{1/4}$,
with a constant $C$ depending only on $\beta$. Using the triangle
inequality and \eqref{eq:sc-floor}, we have
\[
 \|a_N-\alpha_\beta\|_{L^1}
 \le\beta^{-1}\bigl(\|\gamma_N-\gamma^{\rm tr}_\lambda\|_{L^1}
       +\|\gamma^{\rm tr}_\lambda-\gamma_\beta\|_{L^1}\bigr)
 \le C_\beta[\rho^{-3/4}\eta^{1/4}+\lambda^2].
\]
The distance is also at most one, since both functions take values in
$[0,1]$. This proves \eqref{eq:lz-profileconvergence}.

Since $N\ge2$ and $0\le t<1$, we have $d_N\in[\beta^2/2,\beta^2]$ and
$|d_N-\beta^2|=\beta^2t/N\le\beta^2/N$. Then
\cref{lem:scalarcontinuity}, applied with $(a,d)=(a_N,d_N)$ and
$(\widetilde a,\widetilde d)=(\alpha_\beta,\beta^2)$, and
\cref{lem:reference} imply that, for every $q\in[0,1]$, whether or not
$q$ lies in $\supp\mu_\beta$,
\[
 k_{a_N,d_N}(q)\le k_{\alpha_\beta,\beta^2}(q)
 +C_\beta\bigl(\|a_N-\alpha_\beta\|_{L^1}+N^{-1}\bigr)
 \le1+C_\beta\bigl(\|a_N-\alpha_\beta\|_{L^1}+N^{-1}\bigr).
\]
Together with \eqref{eq:lz-profileconvergence}, this yields the bound
on $\varepsilon_N$ in \eqref{eq:lz-stability}. Finally, since
$k_{a_N,d_N}\le1+\varepsilon_N$ pointwise, \cref{lem:Jdom}, applied to
the profile $a_N$ with $\nu=\nu_N$, implies that
$J^{\scal}_{a_N,d_N}\preceq M_{k_{a_N,d_N}}\preceq(1+\varepsilon_N)I$ on
$\mathcal H_{\nu_N}$. Note that only the multipliers of the two
profiles are compared; the operator inequality is proved on the
transport space of the minimizer itself. We choose $C_{\rm st}\ge1$
larger than the constants in the two bounds; it depends only on
$\beta$.
\end{proof}
